\section{Perturbed core costs}\label{app:smoothed}
Random linear perturbations of the core remove this effect in
expectation. For $\sigma>0$ and an integer $M\ge2$ let $\Gamma_M(\sigma)$
be the uniform distribution on the $M$ points
$-\sigma+2\sigma t/(M-1)$, $t=0,\ldots,M-1$.

\begin{lemma}[Finite uniform law]\label{lem:cr-law}
If $\gamma\sim\Gamma_M(\sigma)$ and $I$ is an interval of length
$\ell\ge0$, then $\Pr[\gamma\in I]\le\ell/(2\sigma)+1/M$.
\end{lemma}
\begin{proof}
Consecutive atoms are $2\sigma/(M-1)$ apart, so $I$ contains at most
$\ell(M-1)/(2\sigma)+1$ atoms. Divide by $M$ and use $(M-1)/M\le1$.
\end{proof}

\begin{theorem}[Smoothed query count]\label{thm:cr-smoothed}
Let $F_\gamma(x)=F_0(x)+\sum_{i\in\mathcal C}\gamma_ix_i$, where $F_0$ is
fixed (it may contain residual data drawn independently of $\gamma$)
and the $\gamma_i$ are independent with law $\Gamma_M(\sigma)$. Let $J\ge1$
and $M\ge2^{J-1}$. The core search is correct for every draw, and the
expected number of queries made by levels $0,\ldots,J$ is at most
\[
2^k+8^kJH_k,\qquad
H_k=\Bigl[3+\Bigl(1+\frac k2\Bigr)\frac L{2\sigma}\Bigr]^k .
\]
\end{theorem}
\begin{proof}
Correctness is Theorem~\ref{thm:cr-search}. Write
$V_\gamma(v)=V_0(v)+\gamma^{\T}v$ with
$V_0(v)=\min_zF_0(v,z)$; $V_0$ does not depend on $\gamma$, and by
Lemma~\ref{lem:valuefunction}(a) each $t\mapsto V_0(v+t\mathbf e_i)-Lt^2/2$ is
concave. Fix $j\le J-1$, $h=h_j$, a grid point $v$ and a coordinate $i$
with $v_i\in(0,1)$, so $v\pm h\mathbf e_i\in[0,1]^k$. If
$V_\gamma(v)\le\OPT_\gamma+2e_j$, then
$V_\gamma(v)\le V_\gamma(v\pm h\mathbf e_i)+2e_j$, that is,
\[
\frac{V_0(v)-V_0(v+h\mathbf e_i)-2e_j}{h}\ \le\ \gamma_i\ \le\
\frac{V_0(v-h\mathbf e_i)-V_0(v)+2e_j}{h}.
\]
This interval $I_i(v)$ depends on $V_0$, $v$ and $i$ only. Its length is
$[V_0(v+h\mathbf e_i)+V_0(v-h\mathbf e_i)-2V_0(v)+4e_j]/h\le Lh+4e_j/h
=Lh(1+k/2)$, by concavity of $t\mapsto V_0(v+t\mathbf e_i)-Lt^2/2$ at
$t=-h,0,h$. By independence and Lemma~\ref{lem:cr-law},
\[
\Pr\bigl[V_\gamma(v)\le\OPT_\gamma+2e_j\bigr]\le
\prod_{i:\,v_i\in(0,1)}\Bigl[\frac{Lh(1+k/2)}{2\sigma}+\frac1M\Bigr].
\]
Summing over the $(2^j+1)^k$ grid points coordinate by coordinate, with
two boundary values (factor at most $1$) and $2^j-1$ interior values,
\[
\E N_j\le\Bigl[2+(2^j-1)\Bigl(\frac{Lh_j(1+k/2)}{2\sigma}+\frac1M\Bigr)\Bigr]^k
\le H_k ,
\]
since $(2^j-1)h_j\le1$ and $2^j-1<M$. Lemma~\ref{lem:cr-charge} and
summation over $j=0,\ldots,J-1$ complete the proof.
\end{proof}

\begin{remark}[Meaning of the smoothed count]\label{rem:cr-smoothed}
Correctness never depends on the draw, and the instance being solved is
the perturbed one. The bound is linear in the number $J$ of levels,
whereas Example~\ref{ex:cr-flat} shows that the deterministic count can be
exponential in $J$. The residual dimension and graph enter only the
polynomial cost of each query. Each query has encoding length polynomial
in $I+J+\log M$, so with the cut oracle the expected bit work is
$8^kH_k\poly(I+\log(1/\varepsilon)+\log M)$ with an absolute exponent.
The bound is not a device for approximating the unperturbed problem: the
perturbation can change values by $k\sigma$, so accuracy $\delta$ would
require $\sigma\le\delta/k$, and then $H_k$ grows like $\delta^{-k}$,
worse than the deterministic worst case $(kL/\delta)^{k/2}$. With an
approximate oracle satisfying~\eqref{eq:cr-oracle} with $\eta_j\le ae_j$,
the same proof (applied to the true values $V_\gamma$) gives the bound
with $1+k/2$ replaced by $1+(1+a)k/2$. The noise-confinement argument is
in the spirit of isolation lemmas in smoothed analysis
\cite{BeierVocking2004,RoglinVocking2007}.
\end{remark}

